% Part: incompleteness
% Chapter: representability-in-q
% Section: prim-rec

\documentclass[../../../include/open-logic-section]{subfiles}

\begin{document}

\olfileid{inc}{req}{pri}
\olsection{د بنسټيز بازګښت شبيه کول}

اوس ښودلے شو چې د بنسټيز بازګښت په وسيله تعريف د بېټا تابعې په
کارولو سره د منظم اقل موندلو له لارې «شبيه» کېدے شي۔ فرض کړئ
$f(\vec x)$ او~$g(\vec x, y, z)$ لرو۔ بيا هغه تابعه~$h(\vec x,y)$
چې له~$f$ او~$g$ څخه په بنسټيز بازګښت تعريف شوې وي، دا ده:
\begin{align*}
h(\vec x, 0) & =  f(\vec x) \\
h(\vec x, y+1) & =  g(\vec x, y, h(\vec x, y)).
\end{align*}
\textbf{د ژباړې د نسخې سمون (OLCMP-084):} سرچينې له دې نمايش څخه
مخکښې د تابعې د ارګومېنټونو نااړوند لړ ليکلے و؛ دلته د دواړو
بازګشتي معادلو له ارګومېنټونو سره سم لړ کارول شوے دے۔
بايد وښيو چې~$h$ يوازې د ترکيب او منظم اقل موندلو په کارولو، د
بنسټيزو تابعو او هغو تابعو په مرسته چې له بنسټيزو تابعو څخه د ترکيب
او منظم اقل موندلو له لارې تعريف شوې وي (لکه~$\beta$)، له~$f$ او
$g$ څخه تعريفېدے شي۔

\begin{lem}
\ollabel{lem:prim-rec}
که~$h$ له~$f$ او~$g$ څخه په بنسټيز بازګښت تعريفېدے شي، نو له~$f$،
$g$، تابعو~$\Zero$، $\Succ$، $\Proj{n}{i}$، $\Add$، $\Mult$،
$\Char{=}$ څخه د ترکيب او منظم اقل موندلو په کارولو تعريفېدے شي۔
\end{lem}

\begin{proof}
لومړے يوه مرستندويه تابعه~$\hat h(\vec x, y)$ تعريف کړئ چې تر ټولو
وړوکے عدد~$d$ بېرته ورکوي چې~$d$ داسې لړۍ کوډوي چې لاندې شرطونه پوره
کوي:
\begin{enumerate}
\item $(d)_0 = f(\vec x)$، او
\item د هر~$i < y$ دپاره، $(d)_{i+1} = g(\vec x, i, (d)_i)$،
\end{enumerate}
چېرته چې اوس~$(d)_i$ د~$\beta(d,i)$ لنډون دے۔ په بله وينا،
$\hat h$ دا لړۍ بېرته ورکوي:
$\tuple{h(\vec x, 0), h(\vec x, 1), \dots, h(\vec x, y)}$۔
$\hat h$ داسې ليکلے شو:
\[
\hat h(\vec x, y) = \umin{d}{(\beta(d,0) = f(\vec x) \land \bforall{i <
  y}{\beta(d,i+1) = g(\vec x, i,\beta(d,i)})}.
\]
پام وکړئ: دلته هېڅ بنسټيز بازګښت پکار نۀ دے، يوازې اقل موندل دي۔
هغه تابعه چې اقل ئې مومو د بېټا تابعې د لمې
\olref[bet]{lem:beta} له مخې منظمه ده۔

خو اوس لرو:
\[
h(\vec x, y) = \beta(\hat h(\vec x, y), y),
\]
نو~$h$ له بنسټيزو تابعو څخه يوازې د ترکيب او منظم اقل موندلو په
کارولو تعريفېدے شي۔
\end{proof}

\end{document}
